\chapter{Referencial Teórico}

Este capítulo apresenta os principais conceitos que fundamentam a avaliação da eficácia de suítes de testes automatizados por meio de testes de mutação. Inicialmente, são abordados os fundamentos da Engenharia de Software Experimental e a mineração de repositórios de software, que servem de base para a investigação empírica proposta. Em seguida, são apresentados os conceitos relacionados aos testes automatizados, às métricas de cobertura e à capacidade de detecção de defeitos. Por fim, são discutidos os fundamentos dos testes de mutação, seus operadores, métricas, limitações e sua aplicação no contexto de aplicações web.

\todo[inline, color=pink]{Exemplo de nota ao longo do texto, na cor rosa.}


\section{Engenharia de Software Experimental}

A Engenharia de Software envolve decisões sobre métodos de desenvolvimento, ferramentas, técnicas de programação e práticas de qualidade. Para avaliar essas decisões, não basta considerar apenas os benefícios esperados de determinada abordagem. É necessário investigar seu comportamento em situações reais ou controladas, utilizando evidências que permitam compreender seus resultados.

Segundo \citeonline{wohlin2015}, a investigação empírica é importante na Engenharia de Software porque permite avaliar técnicas, métodos e ferramentas a partir de observações e dados coletados durante sua aplicação. Dessa forma, as decisões podem ser fundamentadas em resultados observáveis, em vez de dependerem exclusivamente de experiências individuais ou pressupostos teóricos.

Os autores apresentam diferentes estratégias de investigação empírica, como questionários, estudos de caso e experimentos controlados. Embora possuam características distintas, essas estratégias não são necessariamente concorrentes, pois a escolha depende do objetivo da pesquisa, do contexto investigado e das condições disponíveis para a coleta e análise dos dados.

Nos experimentos controlados, o pesquisador busca avaliar os efeitos de determinados fatores sobre variáveis de interesse, procurando controlar outras condições que possam influenciar os resultados. Já os estudos de caso investigam fenômenos em seus contextos reais, permitindo compreender aspectos que nem sempre podem ser isolados ou reproduzidos em ambientes controlados \cite{wohlin2015}.

Independentemente da estratégia adotada, a definição clara do problema, das unidades de análise e das métricas utilizadas é necessária para que os resultados possam ser interpretados de maneira adequada.

\todo[inline, color=pink]{Faltou caracterizar melhor a engenharia de software empírica/experimental, os principais métodos e estratégias de pesquisa e relação com a mineração de repositorios de codigo-fonte. Ex: os experimentos, estuddos de caso, etc. Tem todo o livro do wohlin, o guia da revisão. Ainda tem o capitulo 4 que é exatamento o planejamento do estudo emprírico. Ficou muito superficial.}

\subsection{Planejamento e Mensuração em Estudos Empíricos}

\citeonline{wohlin2015} descrevem a experimentação como um processo que abrange a definição e o planejamento do estudo, sua execução, a análise e a interpretação dos dados e a apresentação dos resultados. No planejamento, são definidos o contexto, os objetos de estudo, as variáveis e os procedimentos de coleta e análise. A abordagem \textit{Goal Question Metric} (GQM) orienta a escolha das métricas ao desdobrar os objetivos da investigação em perguntas e selecionar as medidas necessárias para respondê-las.

O planejamento também deve considerar as ameaças à validade interna, externa, de construção e de conclusão, relacionadas, respectivamente, às relações observadas, à generalização dos resultados, à adequação das medidas e à consistência das conclusões \cite{wohlin2015}. Em estudos com diferentes projetos de software, características como linguagem, arquitetura e tamanho das suítes de testes podem influenciar os resultados e devem ser consideradas na sua interpretação.

\subsection{Mineração de Repositórios de Software}

A mineração de repositórios de software, conhecida como \textit{Mining Software Repositories (MSR)} ou \textit{ Repository Mining}, é uma abordagem de pesquisa empírica que utiliza dados armazenados em repositórios de software para investigar características dos sistemas, seu desenvolvimento e sua evolução. Segundo \citeonline{hassan2008}, seu objetivo é analisar e relacionar informações disponíveis nesses repositórios para identificar padrões e produzir conhecimentos que possam apoiar decisões durante o desenvolvimento e a manutenção de software.

Os repositórios armazenam diferentes tipos de informações, como código-fonte, histórico de modificações, registros de defeitos, mensagens de commits e comunicações entre desenvolvedores. \citeonline{hassan2008} classifica essas fontes em repositórios históricos, que registram a evolução dos projetos; repositórios de execução, que armazenam informações sobre o funcionamento dos sistemas; e repositórios de código, que disponibilizam implementações de software. A análise conjunta dessas informações permite investigar aspectos que nem sempre podem ser identificados apenas pela observação do código-fonte.

\citeonline{spadini2018} também destacam que os repositórios contêm informações valiosas sobre o desenvolvimento de sistemas e que sua mineração permite compreender práticas de manutenção, evolução do código e atividades dos desenvolvedores. A partir desses dados, é possível investigar diferentes questões da Engenharia de Software, como identificação de componentes propensos a defeitos, análise de alterações e práticas de teste.

\citeonline{hassan2008} ressalta que uma das principais contribuições da mineração de repositórios é transformar registros utilizados originalmente para armazenar informações em fontes de conhecimento para a tomada de decisões. Por exemplo, o histórico de modificações e os registros de defeitos podem ser relacionados para identificar componentes que apresentaram problemas recorrentes e, assim, apoiar a priorização de atividades de manutenção e teste.

No contexto deste trabalho, a mineração de repositórios será utilizada para selecionar e obter dados de aplicações web de código aberto, especialmente seus códigos-fonte e suítes de testes automatizados. Esses artefatos servirão de base para a aplicação dos testes de mutação, permitindo avaliar a capacidade das suítes em identificar alterações artificiais introduzidas no código.

\subsubsection{Seleção e Extração de Dados}

Na MSR, os objetivos e as questões de pesquisa orientam a escolha das fontes de dados e a definição dos critérios de inclusão e exclusão. \citeonline{vidoni2022} distingue as plataformas que hospedam os projetos, os repositórios selecionados e os dados extraídos de cada um, como código-fonte, histórico de commits e registros de problemas.

Os critérios de seleção devem considerar características relevantes para a pesquisa e possíveis fontes de viés, como repositórios duplicados ou derivados e projetos inativos \cite{vidoni2022}. A extração consiste em recuperar os dados necessários à análise e pode ser automatizada com ferramentas como o PyDriller, que permite obter commits e informações sobre arquivos modificados em repositórios Git \cite{spadini2018}.

\subsubsection{Análise dos Dados e Reprodutibilidade}

Após a extração, os dados devem ser examinados e preparados para análise. Registros incompletos, inconsistências e diferenças entre repositórios podem afetar a interpretação dos resultados \cite{hassan2008}. Por isso, é importante explicitar como esses problemas foram tratados e discutir as ameaças à validade associadas à seleção, à coleta e à análise \cite{vidoni2022}.

Para favorecer a reprodutibilidade, devem ser registrados os repositórios e as versões analisadas, as datas de coleta, os dados extraídos, as ferramentas utilizadas e os procedimentos de processamento e análise. Esse registro permite rastrear a origem dos resultados e repetir o procedimento sobre as mesmas versões dos repositórios \cite{vidoni2022}.


\todo[inline, color=pink]{Incluir uma subseção que  apresente o assunto de qualidade de produto de software, relacionando modelos de qualidade de referencia, como a ISO/IEC 25010, com a avaliação de testes automatizados. Destaquem que o teste de software é uma atividade de verificação e validação que contribui para a avaliação da qualidade do produto.}

\section{Testes de Software}

Os testes de software são atividades utilizadas para avaliar o comportamento de um sistema e identificar diferenças entre os resultados observados e aqueles esperados. Segundo \citeonline{ammann2016}, os testes permitem verificar o comportamento de um programa por meio de sua execução, contribuindo para a identificação de defeitos e para a redução dos riscos associados à utilização do software.

Entretanto, a execução de testes não garante a ausência de defeitos. Conforme discutido por \citeonline{ammann2016}, não é possível, na maioria dos sistemas, verificar todas as combinações de entradas e condições de execução. Dessa forma, os testes são selecionados com base em critérios que procuram avaliar os comportamentos mais relevantes do programa.

A capacidade de uma suíte identificar problemas depende de fatores como a escolha dos casos de teste, os dados utilizados, os comportamentos exercitados e as verificações realizadas sobre os resultados. Por esse motivo, além da execução dos testes, é necessário avaliar sua adequação e capacidade de detectar comportamentos incorretos \cite{ammann2016}.

\subsection{Testes Automatizados e Suítes de Testes}

A automação de testes consiste na utilização de ferramentas e scripts para executar verificações de software com menor necessidade de intervenção manual. De acordo com \citeonline{ammann2016}, a automação permite organizar a execução dos testes e realizar verificações de maneira sistemática, facilitando sua repetição ao longo do desenvolvimento.

Uma suíte de testes corresponde a um conjunto organizado de casos de teste utilizados para avaliar determinados comportamentos de um sistema. Esses casos podem conter entradas, condições de execução e verificações dos resultados obtidos. \citeonline{ammann2016} destacam a importância de estruturar os casos de teste de maneira que sua execução e avaliação possam ser realizadas de forma automatizada.

As suítes automatizadas podem atuar em diferentes níveis de teste. Os testes unitários verificam componentes individuais, como funções e métodos. Os testes de integração avaliam a interação entre componentes ou serviços, enquanto os testes de sistema procuram verificar o comportamento da aplicação como um todo. Os testes de ponta a ponta (\textit{end-to-end}) também podem ser utilizados para validar fluxos completos de utilização do sistema \cite{ammann2016}.

A automação também está relacionada às práticas de integração contínua, nas quais verificações são executadas frequentemente durante o desenvolvimento. \citeonline{ammann2016} abordam a integração contínua como uma prática que permite executar testes regularmente, contribuindo para identificar problemas após modificações no código-fonte.

É importante salientar que executar testes frequentemente não significa necessariamente que eles sejam capazes de identificar todos os defeitos introduzidos no software. Por isso, além de verificar se uma suíte executa corretamente, é necessário avaliar sua capacidade de reconhecer comportamentos incorretos.

\subsection{Eficácia de Suítes de Testes}

A eficácia de uma suíte de testes está relacionada à sua capacidade de identificar defeitos no software avaliado. Segundo \citeonline{just2014}, uma boa suíte de testes é aquela que consegue detectar defeitos reais. Entretanto, medir diretamente essa capacidade é difícil, pois o conjunto de defeitos existentes em um programa geralmente não é conhecido.

Diante dessa dificuldade, pesquisadores utilizam critérios indiretos para avaliar a adequação dos testes. \citeonline{ammann2016} apresentam diferentes critérios de cobertura que podem orientar a elaboração e a avaliação de casos de teste. Esses critérios estabelecem requisitos que uma suíte deve satisfazer, como a execução de determinadas instruções, caminhos ou condições do programa.

Entretanto, satisfazer um critério de cobertura não significa necessariamente possuir maior capacidade de detecção de defeitos. \citeonline{inozemtseva2014} investigaram a relação entre tamanho, cobertura e eficácia de suítes de testes, utilizando aproximadamente 31 mil suítes geradas para cinco sistemas Java. Os autores identificaram uma correlação baixa a moderada entre cobertura e eficácia quando o tamanho das suítes era controlado.

A avaliação da eficácia deve considerar que diferentes suítes podem apresentar resultados distintos, mesmo quando executam partes semelhantes do programa. Por exemplo, dois conjuntos de testes podem executar uma mesma função, mas apenas um deles verificar corretamente o resultado produzido. Nesse caso, a capacidade de identificar um defeito poderá ser diferente, mesmo que ambos apresentem cobertura semelhante.

Assim, a quantidade de testes ou a proporção de código executada não são suficientes, isoladamente, para determinar a eficácia de uma suíte. Conforme discutido por \citeonline{just2014}, a utilização de defeitos artificiais, denominados mutantes, constitui uma alternativa para avaliar a capacidade dos testes de identificar comportamentos incorretos quando não existe um conjunto conhecido de defeitos reais.

\subsection{Oráculos de Teste}

Segundo \citeonline{barr2015}, um oráculo de teste (\textit{test oracle}) é um mecanismo que determina se o comportamento observado durante a execução de um teste corresponde ao esperado. Essa avaliação pode se basear em especificações, modelos, propriedades do sistema ou resultados previamente conhecidos.

Em testes automatizados, as asserções permitem implementar essas verificações ao estabelecer condições que devem ser satisfeitas. Por exemplo, um teste de uma função de desconto pode comparar o valor retornado com o valor esperado. Se verificar apenas se a função terminou sem erros, o teste pode passar mesmo quando o desconto calculado está incorreto \cite{ammann2016,barr2015}.

Nos testes de mutação, a capacidade de detectar um mutante depende do que o oráculo verifica. Uma alteração pode ser executada e produzir um efeito observável, mas o mutante pode sobreviver se os testes não verificarem o resultado afetado.

\subsection{Testes em Aplicações Web}

Em aplicações web, uma funcionalidade pode combinar o processamento de dados no servidor e a atualização da interface no navegador. Os testes automatizados podem verificar o comportamento em diferentes pontos, como a resposta de uma API ou o conteúdo exibido ao usuário. Enquanto um teste de API pode conferir os dados retornados pelo servidor, um teste de ponta a ponta pode verificar o resultado de um fluxo completo pela interface \cite{li2014webtesting}.

\citeonline{li2014webtesting} destacam que a interface pode ser atualizada de forma assíncrona, sem recarregar a página. Nesses casos, o teste precisa considerar o momento em que o resultado esperado fica disponível para verificação. Ao enviar um formulário, por exemplo, o teste pode aguardar que a mensagem de confirmação apareça antes de conferir seu conteúdo.

\section{Testes de Mutação}

Os testes de mutação (\textit{Mutation Testing}) constituem uma técnica utilizada para avaliar a adequação de suítes de testes por meio da introdução de alterações artificiais em um programa. Segundo \citeonline{ammann2016}, a técnica faz parte dos critérios de teste baseados em sintaxe, nos quais regras de transformação são utilizadas para produzir variações de um artefato de software.

Segundo \citeonline{just2014}, a utilização de defeitos artificiais permite avaliar a capacidade de detecção de defeitos de uma suíte, especialmente quando não existe um conjunto conhecido de defeitos reais que possa ser utilizado como referência. Dessa forma, os testes de mutação também podem contribuir para identificar limitações nas suítes, orientar a elaboração de novos casos de teste e comparar diferentes estratégias de teste.

\subsection{Fundamentos dos Testes de Mutação}

A aplicação dos testes de mutação parte de um programa original, sobre o qual são realizadas pequenas alterações sintáticas. Essas modificações são produzidas por operadores de mutação, definidos de acordo com as características da linguagem de programação utilizada.

\citeonline{ammann2016} definem um operador de mutação como uma regra que especifica possíveis variações sintáticas de um artefato. O resultado da aplicação de um operador é denominado mutante. No caso da mutação de programas, cada mutante representa uma versão do código-fonte que contém uma alteração em relação à implementação original.

Os autores destacam que os mutantes não são, por si só, casos de teste. Sua finalidade é estabelecer requisitos que os testes devem satisfazer. Assim, o objetivo imediato da técnica é identificar casos de teste capazes de distinguir o programa original de suas versões modificadas, o que estabelece o critério de cobertura por mutação (Mutation Coverage).

A fundamentação dos testes de mutação também está relacionada a dois pressupostos teóricos: a hipótese do programador competente (\textit{Competent Programmer Hypothesis}) e o efeito de acoplamento (\textit{Coupling Effect}).

A hipótese do programador competente considera que desenvolvedores geralmente produzem programas próximos do comportamento correto, de modo que muitos defeitos podem ser representados por pequenas alterações na implementação. Já o efeito de acoplamento estabelece que testes capazes de identificar defeitos simples também podem detectar defeitos mais complexos relacionados a eles \cite{ammann2016}.

Esses pressupostos ajudam a justificar a utilização de alterações artificiais como instrumento de avaliação. Entretanto, não significam que todos os defeitos reais possam ser representados pelos operadores disponíveis \cite{just2014}.

\subsection{Processo de Aplicação dos Testes de Mutação}

A aplicação dos testes de mutação envolve a geração de versões modificadas de um programa e a execução de testes para verificar se essas alterações podem ser detectadas.

\citeonline{ammann2016} apresentam um processo que começa com a seleção do programa a ser testado e a geração dos mutantes. Em seguida, podem ser utilizadas técnicas para identificar e eliminar mutantes equivalentes. Os casos de teste são executados inicialmente sobre o programa original e, posteriormente, sobre as versões modificadas.

Quando um caso de teste produz um resultado observável diferente entre o programa original e um mutante, considera-se que esse mutante foi \textbf{morto} (\textit{killed}). Após sua detecção, não é necessário continuar executando outros testes sobre ele para satisfazer o requisito de mutação correspondente.

Por outro lado, um mutante é classificado como \textbf{sobrevivente} (\textit{survived}) quando nenhum dos testes executados consegue distinguir seu comportamento daquele apresentado pelo programa original. A sobrevivência indica que a alteração não foi detectada pela suíte nas condições avaliadas, mas não permite concluir, isoladamente, que os testes sejam inadequados.

Além disso, alguns mutantes podem apresentar comportamento equivalente ao programa original. Nesses casos, a sobrevivência não representa necessariamente uma deficiência da suíte, pois não existe uma diferença funcional a ser identificada.

\citeonline{ammann2016} distinguem alguns tipos de mutantes que precisam ser considerados durante a avaliação:

\begin{itemize}
    \item \textbf{Mutantes inválidos (\textit{stillborn mutants}):} apresentam alterações sintaticamente incorretas que impedem a compilação ou interpretação do programa. Esses mutantes devem ser evitados durante a geração ou descartados da avaliação;

    \item \textbf{Mutantes equivalentes (\textit{equivalent mutants}):} possuem alterações sintáticas, mas apresentam o mesmo comportamento funcional do programa original. Como não produzem diferenças observáveis, não podem ser mortos por casos de teste;

    \item \textbf{Mutantes triviais (\textit{trivial mutants}):} apresentam alterações facilmente detectáveis, podendo ser mortos por praticamente qualquer caso de teste que exercite adequadamente a região modificada.
\end{itemize}

A identificação de mutantes equivalentes constitui uma dificuldade particular. Segundo \citeonline{ammann2016}, determinar a equivalência entre dois programas é, no caso geral, um problema indecidível e  não existe um procedimento automático capaz de resolver todos os casos possíveis.

Após a execução dos testes, os resultados são utilizados para calcular a proporção de mutantes mortos e avaliar a adequação da suíte. O processo apresentado por \citeonline{ammann2016} também prevê a possibilidade de criar novos casos de teste direcionados aos mutantes sobreviventes, repetindo a avaliação até atingir um nível de adequação considerado satisfatório.

No contexto deste trabalho, os fundamentos desse processo serão utilizados para avaliar suítes automatizadas já existentes em projetos de código aberto. Dessa forma, a execução dos mutantes permitirá identificar quais alterações são detectadas pelos testes disponíveis, sem exigir necessariamente a criação de novos casos de teste.

\subsubsection{Mutação Forte e Mutação Fraca}

Os testes de mutação podem ser classificados de acordo com o momento em que as diferenças entre o programa original e o mutante são verificadas.

\citeonline{ammann2016} explicam essa distinção a partir de três condições necessárias para a manifestação de uma falha: alcance (\textit{Reachability}), infecção (\textit{Infection}) e propagação (\textit{Propagation}).

O alcance ocorre quando um caso de teste executa a região do código na qual a alteração foi introduzida. A infecção ocorre quando essa alteração produz uma diferença no estado interno do programa. Por fim, a propagação ocorre quando essa diferença alcança um resultado observável da execução.

Na \textbf{mutação forte} (\textit{Strong Mutation}), um mutante é considerado morto quando um caso de teste produz um resultado observável diferente daquele apresentado pelo programa original. Dessa forma, a detecção depende das condições de alcance, infecção e propagação.

Já na \textbf{mutação fraca} (\textit{Weak Mutation}), a avaliação ocorre imediatamente após a execução do trecho modificado. Nesse caso, o mutante pode ser considerado morto quando a alteração produz uma diferença no estado interno do programa, mesmo que essa diferença não seja propagada até a saída final.

Portanto, um mutante pode ser morto segundo o critério de mutação fraca e permanecer sobrevivente segundo o critério de mutação forte. Essa distinção é importante porque demonstra que a execução de uma alteração e a modificação do estado interno não garantem, necessariamente, que seu efeito seja identificado nos resultados do programa \cite{ammann2016}.

\subsection{Operadores de Mutação e Tipos de Mutantes}

Os operadores de mutação são responsáveis por definir as alterações que serão introduzidas no programa original. Sua escolha é importante porque determina os tipos de modificações que poderão ser avaliadas.

Segundo \citeonline{ammann2016}, os operadores são geralmente desenvolvidos com dois objetivos: representar erros comuns cometidos durante a programação ou incentivar a elaboração de casos de teste que satisfaçam determinados critérios de adequação.

Os autores apresentam diferentes operadores aplicáveis ao código-fonte, entre eles aqueles que modificam expressões aritméticas, relacionais, condicionais, atribuições e referências a variáveis.

O \textbf{operador de substituição aritmética} (\textit{Arithmetic Operator Replacement} -- AOR), por exemplo, modifica operadores como soma, subtração, multiplicação e divisão. Assim, uma expressão como \texttt{x = a + b} pode ser modificada para \texttt{x = a - b}, produzindo um comportamento diferente para determinadas entradas.

O \textbf{operador de substituição relacional} (\textit{Relational Operator Replacement} -- ROR) altera operadores utilizados em comparações, como \texttt{>}, \texttt{<}, \texttt{>=} e \texttt{<=}. Essas modificações permitem avaliar se os testes conseguem identificar alterações nas condições de decisão do programa.

Outro exemplo é o \textbf{operador de substituição condicional} (\textit{Conditional Operator Replacement} -- COR), que modifica conectores lógicos utilizados nas expressões condicionais. Uma condição que utiliza o operador \texttt{\&\&}, por exemplo, pode ser modificada para utilizar \texttt{||}, alterando as condições necessárias para que determinada expressão seja considerada verdadeira.

Além desses, \citeonline{ammann2016} apresentam operadores de inserção e remoção de operadores unários, substituição de operadores de atribuição e substituição de referências a variáveis.

Existem também operadores desenvolvidos para exigir determinados comportamentos dos testes. Um exemplo é o operador \textit{Bomb Statement Replacement} (BSR), que substitui uma instrução por uma operação que provoca uma falha quando executada. Dessa forma, o mutante somente pode ser morto por um teste que alcance a instrução modificada, estabelecendo uma relação com o critério de cobertura de instruções.

Os exemplos demonstram que os operadores não possuem necessariamente a mesma finalidade. Enquanto alguns representam alterações semelhantes a possíveis erros de programação, outros são utilizados para exigir que determinados elementos do código sejam exercitados ou verificados pelos testes.

Os mutantes também podem ser classificados de acordo com a quantidade de alterações introduzidas. Os \textbf{mutantes de primeira ordem} (\textit{First-Order Mutants} -- FOM) são produzidos pela aplicação de uma única alteração, enquanto os \textbf{mutantes de ordem superior} (\textit{Higher-Order Mutants} -- HOM) combinam duas ou mais alterações em uma mesma versão do programa.

\citeonline{ammann2016} destacam que a aplicação de uma alteração por mutante é geralmente preferida na mutação de programas, embora existam situações em que combinações de alterações possam ser úteis.

A escolha e a quantidade dos operadores também influenciam o custo da análise. Dependendo das características do programa, sua aplicação pode produzir um grande número de mutantes, exigindo diversas execuções da suíte de testes.

Nesse contexto, \citeonline{ammann2016} apresentam duas abordagens gerais para reduzir a quantidade de mutantes avaliados: a \textbf{amostragem aleatória}, que seleciona parte dos mutantes gerados, e a \textbf{mutação seletiva} (\textit{Selective Mutation}), que utiliza um conjunto reduzido de operadores considerados eficazes.

Dessa forma, a definição dos operadores deve considerar a linguagem de programação, as características do sistema e os objetivos da avaliação.

\subsection{Indicador de Mutação e Avaliação da Eficácia}

O Indicador de Mutação (\textit{Mutation Score}) é uma métrica utilizada para quantificar a adequação de uma suíte de testes em relação ao conjunto de mutantes gerados.

Segundo \citeonline{ammann2016}, o critério de cobertura por mutação considera que cada mutante estabelece um requisito de teste, que consiste em encontrar ao menos um caso capaz de identificá-lo. Dessa forma, o Indicador de Mutação representa a proporção de mutantes mortos em relação ao total de mutantes não equivalentes considerados na avaliação.

O cálculo da métrica pode ser representado pela seguinte expressão:

\begin{equation}
    MS = \frac{K}{M - E} \times 100
    \label{eq:mutation-score}
\end{equation}

Em que $K$ representa a quantidade de mutantes mortos, $M$ corresponde ao total de mutantes válidos gerados e $E$ representa a quantidade de mutantes equivalentes entre esses mutantes.

Um resultado de 90\%, por exemplo, indica que a suíte conseguiu detectar 80\% dos mutantes válidos e não equivalentes considerados na avaliação.

Entretanto, alcançar resultado alto como esse pode ser difícil, na prática. Por esse motivo, os autores também apresentam a possibilidade de estabelecer um limiar mínimo de adequação, utilizado para orientar a continuidade do processo de testes.

É importante destacar que a adequação é relativa ao conjunto de mutantes utilizado. Uma suíte que apresenta um Indicador de Mutação elevado não necessariamente consegue detectar todos os defeitos reais do programa, pois os mutantes representam apenas as alterações produzidas pelos operadores selecionados.

A relação entre detecção de mutantes e defeitos reais também é importante para a interpretação dessa métrica. \citeonline{just2014} discutem a utilização de mutantes como uma aproximação para avaliar a capacidade de detecção de defeitos reais, embora reconheçam que nem todos os defeitos possam ser representados pelas alterações artificiais disponíveis.

Dessa forma, sua interpretação pode ser complementada por informações como cobertura de código, quantidade de mutantes sobreviventes, operadores aplicados e características dos casos de teste avaliados. Esses indicadores permitem compreender melhor os resultados da avaliação e identificar possíveis limitações das suítes automatizadas.


\todo[inline, color=pink]{Essa parte de testes automatizados e testes de mutação ficou muito boa. Contudo, como você precisaram acertar a parte de engenharia de software experimental, talvez seja necessário resumir mais um pouco. Vejam essa equilíbrio no total de páginas do tcc. }
